\subsection{All-place evaluation and harmonic quartic initialization}

\begin{statement}{A.112}{Proved}
{The harmonic quartic joint characters give no finite initialization obstruction}
Use the curves, maps and normalized homomorphism
$\Theta_R$ of A.108.
For every $Q\in E(\Q)$ and every odd prime $p$,
the actual localized class is zero:
\[
       \bar r^{\,*}[\kappa_Q]=0
                       \quad\text{in }H^1(\Q_p,J_4).
\tag{A.112.1}
\]
The same equality holds over $\R$.
For every negative-triple source (A.108.5), the
quaternion (A.108.6) therefore evaluates to zero
at every local point of $C_4$, at every place.
At two this is the evaluation theorem A.111;
no zero class in $H^1(\Q_2,J_4)$ is asserted here.

Let $c\in\Q\setminus\{0,1,\tfrac12\}$ be any new
puncture. Fix a finite set of places and one
ordinary point of $C_4$ at each, with $T\ne c$.
There exists $a_c\in\Q^\times$ such that the
augmented quartic cover
\[
 \left\{
 \begin{array}{l}
 T-c_i=(R-c_i)X_i^4,\quad c_i=0,1,\tfrac12,\\
 T-c=a_c Z^4
 \end{array}
 \right\}
 \longrightarrow C_4,\qquad
             (T,X_i,Z)\longmapsto(T,X_i)
\tag{A.112.2}
\]
has an ordinary point over every completion and
lifts each prescribed old point.
This concerns one finite quartic cover and a
finite prescribed set of old points.

The source norm obstruction can nevertheless
be nonzero before the map in (A.112.1).
For
\[
 \begin{gathered}
 d=138,\quad R=24/23,\quad P_R=(150,720),\\
 Q=(162,1080),\quad c=-2/23,\quad \rho=3/2,
 \end{gathered}
\tag{A.112.3}
\]
one has $[\kappa_Q]\ne0$ in $H^1(\Q_3,E)$,
whereas its image in $H^1(\Q_3,J_4)$ is zero.
Thus this cohomology map has a concrete nonzero
kernel element, despite the injectivity of the
Jacobian homomorphism in A.108.
\end{statement}
\begin{proof}
\emph{Good reduction and the actual unramified norm.}
If $p\equiv1\pmod4$, then $i\in\Q_p$ and the
cocycle (A.108.3) is already zero.
Suppose $p\equiv3\pmod4$; then
$K=\Q_p(i)$ is the unramified quadratic extension.

If $v_p(d)$ is even, rescale $x,y$ by
$p^{2a},p^{3a}$, where $2a=v_p(d)$.
This is an actual $\Q_p$-isomorphism to the model
with unit $d'=d/p^{2a}$.
Its discriminant $64d'^6$ is a unit, so it has
good reduction. The norm
\[
                 N:E(K)\longrightarrow E(\Q_p)
\tag{A.112.4}
\]
is onto, as follows directly from reduction.
Let $\varphi$ be Frobenius on the elliptic special
fibre over $\F_p$.
The endomorphism $1+\varphi$ is a nonzero isogeny.
For $\bar Q\in E(\F_p)$ choose a geometric
preimage $\bar S$ under it.
Then
$(\varphi^2-1)\bar S=(\varphi-1)\bar Q=0$,
so $\bar S\in E(\F_{p^2})$ and its norm is $\bar Q$.
Lift it to $S\in E(K)$ by smooth Hensel lifting.
The error $Q-N(S)$ lies in the formal kernel
over $\Q_p$.
Multiplication by two is an automorphism of
that formal group, since its linear coefficient
is the unit two.
Its half $A\in E(\Q_p)$, considered over $K$,
has norm $2A$.
Thus $S+A$ has norm $Q$, proving (A.112.4).
The source norm criterion of A.108 gives
$[\kappa_Q]=0$.

\emph{Odd valuation and the exact component representatives.}
If $v_p(d)$ is odd, the same rescaling puts
$v_p(d)=1$.
Write $d=p\delta$, with $\delta$ a unit.
The special cubic is $y^2=x^3$.
Its nonsingular locus, with infinity as identity,
is the additive group.
Let $E_0(\Q_p)$ be the points reducing to this
nonsingular locus.
Its reduction is $\mathbb G_a(\F_p)$ and its
kernel is the one-dimensional formal group.
Multiplication by two is bijective on both.
Lifting a half on the smooth special fibre
and then dividing the formal error proves that
$E_0(\Q_p)$ is uniquely two-divisible.
This uses the actual nonsingular reduction,
without a component-group classification.

Every point is a rational two-torsion translate
of a point in $E_0(\Q_p)$.
The identity and two-torsion points can be treated
directly, so take another point $P=(x,y)$.
If $v_p(x)\le0$, it already has nonsingular
reduction, in the formal neighbourhood of infinity
when the valuation is negative.
For $v_p(x)\ge2$, translation by $T_0=(0,0)$ gives
\[
 x(P+T_0)=-\frac{d^2}{x(P)},\qquad
                v_p(x(P+T_0))=2-v_p(x(P))\le0.
\tag{A.112.5}
\]
For $v_p(x)=1$, put $x=pz$.
The equation $y^2=p^3z(z^2-\delta^2)$ forces
$v_p(z^2-\delta^2)$ to be odd and positive.
Since $p$ is odd, exactly one of $z-\delta,z+\delta$
is divisible by $p$.
Choose $\varepsilon\in\{1,-1\}$ accordingly.
Then $v_p(x-\varepsilon d)\ge2$ and translation by
$T_\varepsilon=(\varepsilon d,0)$ gives
\[
       x(P+T_\varepsilon)
         =\varepsilon d+\frac{2d^2}{x-\varepsilon d}.
\tag{A.112.6}
\]
The second summand has valuation at most zero,
whereas the first has valuation one.
Their sum has valuation at most zero and belongs
to $E_0(\Q_p)$.
Consequently
\[
                    E(\Q_p)=E[2](\Q_p)+2E(\Q_p).
\tag{A.112.7}
\]
The map $Q\mapsto[\kappa_Q]$ has two-torsion image.
For each two-torsion point $T$, choose the branch
class $L_k$ with $q(B_k)=[i]T$.
Equation (A.108.21) makes $-jL_k$ the explicit
primitive of $h\kappa_T$.
Thus (A.112.7) proves (A.112.1).
The rescaling isomorphisms commute with $[i]$
and preserve the corresponding named two-torsion
points, so this is the original map of A.108.

\emph{The real place and the local evaluations.}
The cubic has three distinct real roots.
The real group has two connected components;
its identity component is a circle, on which
multiplication by two is onto.
The point $(0,0)$ lies on the other, bounded
component. Hence
$E(\R)=E[2](\R)+2E(\R)$, and the same primitives
prove $h[\kappa_Q]=0$ over $\R$.

At every odd or real place, (A.108.6) has zero
Hochschild--Serre image and is zero at the old
point. The split comparison in A.88 makes it
the zero Brauer class on the local curve.
Every evaluation is therefore zero.
At two, A.111 gives zero at every ordinary
point with $T\ne c$.
The class is unramified on the proper curve.
Evaluation is locally constant, and the finitely
many excluded fibres are limits of ordinary points
on this smooth local curve.
The same value follows at any local points of those
fibres. This proves the all-place evaluation claim,
without turning A.111 into a cohomology vanishing
statement.

\emph{Every common quartic character and the initialization map.}
For the new puncture put $D_i=(c-c_i)/(R-c_i)$.
A.96 gives the exact common abelian fibre field
\[
 B_4=\Q\left(\sqrt{D_i},
       \sqrt{2r_J}:
       \prod_{i\in J}D_i=-r_J^2,\quad r_J>0\right).
\tag{A.112.8}
\]
Every character is a sum of these quadratic
characters.
For the individual $D_i$, the norm or Steinberg
identity of A.96 makes the cyclic class a global
constant on the old cover.
For a negative singleton or pair, its explicit
quartic norm identities and actual affine coordinate
comparison give the same conclusion.
These identities use only the indicated one or two
old coordinates, so they pull back along the
actual forgetful morphism from $C_4$.
After division by $R-c$, every such class is
zero at the old point, and has zero evaluation
at every local point.

The remaining generator is the negative triple,
whose relation is exactly $F(c)/F(R)=-r^2$.
Its normalized symbol is (A.108.6), with zero
evaluation everywhere by the preceding proof
and A.111.
Thus every character in (A.112.8) has zero
normalized evaluation on every local old point.
For the unnormalized symbols the evaluations are
those of the single global constant $(R-c,\chi)$;
their local invariants sum to zero.

Apply the exact finite-character initialization
theorem A.85 to this old cover and the prescribed
finite collection of old points.
The character obstruction just computed vanishes.
Its construction supplies a rational coefficient
$a_c$ and ordinary points of (A.112.2) everywhere,
while leaving the prescribed old points fixed;
the real sign condition is retained.
If $c=R$, the fibre algebra has the rational root
tuple $(1,1,1)$.
Its common abelian field is then $\Q$, so there
are no character conditions and A.85 applies
directly.
This proves the stated finite initialization.
It neither prescribes an infinite list of local
points nor constructs a cofinal system of old
covers.

\emph{A nonzero source class killed by this cohomology map.}
The two points in (A.112.3) satisfy $E_{138}$.
The source norm quartic of A.108 is
\[
                 z^2=-m^4+972m^2+8640m+2556.
\tag{A.112.9}
\]
For $m\in\Z_3$ its right side is $18-m^4$ modulo
twenty-seven.
If $m$ is a unit, its residue modulo three is two.
If $3\mid m$, its residue modulo twenty-seven is
eighteen, with valuation two and nonsquare unit part.
Neither is a square.
For $v_3(m)<0$, division by $m^4$ leaves a unit
congruent to $-1$ modulo three, also nonsquare.
The points at infinity require $i\in\Q_3$.
The projective norm torsor therefore has no
$\Q_3$-point, so A.108 proves $[\kappa_Q]\ne0$.
Since $v_3(138)=1$, the proved odd-valuation
argument makes its image under $\bar r^{\,*}$ zero.
This supplies the announced cohomology kernel
element. Injectivity of $\bar r^{\,*}$ on abelian
varieties was proved from the generating inertia
groups of $C_4\to C_2$ at the end of A.108.
\end{proof}

These results remove the harmonic quartic
joint-character evaluations from the finite
initialization obstruction.
They assert neither global vanishing of
$\Theta_R(Q)$ nor zero localization at two in
$H^1(\Q_2,J_4)$.
The actual global and local maps remain
(A.108.4), (A.108.16), and (A.112.1).

The exact certificate
\path{code/harmonic_local_initialization_certificate.py},
with retained output, checks the two-torsion translation
formulas on locally soluble odd-place samples, computes
actual finite-field quadratic norms, and verifies every
residue class in the explicit modulo-twenty-seven
source obstruction.

